\section{Known unknowns}\label{sec:upunknowns}

The items below are listed in order of how much they could move the answer.
Each is an \emph{assumption} or an unresolved discrepancy, not a measurement.

\begin{enumerate}
\item \textbf{Live-intermediate fraction} (\smLiveFrac). It dominates on-floor
      storage and has not been measured.
\item \textbf{DR1 compute basis}: \vthirty{} (\smDRoneNodeDays{} node-days)
      against the earlier \wfortyone{} estimate (about 5{,}700). DP2's measured
      cost per visit is closer to the earlier figure (\secref{sec:upcompute}).
      The difference is worth about a factor of two in DRP compute.
\item \textbf{Idle and inter-stage multiplier} (\smIdle). Only the task-level
      ratio of wall time to CPU time (about 1.14) has been measured.
\item \textbf{Release cadence}. This update assumes one data release a year,
      and a release every two or three years is under discussion. Fewer
      releases would each be larger, with fewer campaigns in total. Whether
      that saves hardware depends on whether a campaign is stretched out or
      run more intensively, which has not been decided. Modelling it needs a
      release-schedule parameter, which is not yet implemented.
\item \textbf{Qserv catalogue basis}. The modelled DR1 catalogue is more than
      forty times the measured DP2 catalogue, for about seven times the visits.
      Per-table row counts for DP2 would settle this. They would also show:
      \begin{itemize}
      \item which tables carry overlap (worth up to about 13\% at DR9);
      \item whether ForcedSource and ForcedSourceOnDIAObject share the 41-byte
            row size;
      \item where the independent DR1 estimate of
            \smQsPlanLowPB--\smQsPlanHighPB\,PB comes from.
      \end{itemize}
\item \textbf{Which Qserv deployment was measured}, and whether the integration
      cluster is representative of production. The measured catalogue used
      two-fold replication; production is to use three-fold.
\item \textbf{When the empty NVMe slots are populated}. This decides whether
      Qserv has a node gap at DR1 (\secref{sec:upqserv}).
\item \textbf{Partner shares and capacity}: UKDF's year-by-year profile;
      France's share and capacity; and a review of partner provisioning
      against the current compute basis (\secref{sec:upsplit}).
\item \textbf{Torino speedup} (\smSpeedup). Measured on one generation of
      tasks.
\item \textbf{Tape}: how much shared tape capacity is actually available to
      Rubin, and how much archive is already on the floor.
\item \textbf{Kubernetes allocatable cores}: 48 in planning documents against
      64 hardware threads deployed (DM-54218).
\item \textbf{DR1-scale sizes} for Scarlet models and HiPS maps.
\item \textbf{Hosting and overhead rates}, which were not re-quoted, and a
      price for the fast (NVMe) tier.
\end{enumerate}
